% =====================================================================
%  Report 5
%  Two observations on the Metadata Guideline and the EarthPortal schema
%  Build: latexmk -pdf emmanuel_dataTerra_Report5-Cheatsheet-vs-EarthPortal.tex
%  Data: the metadata schema of EarthPortal (data.earthportal.eu/submission_metadata and
%  /ontology_metadata) and the latest submission of the 68 artefacts, all read on 13 October 2026;
%  the comparison itself comes from notebook 04 of github.com/EarthPortal/SA_Management.
% =====================================================================
\documentclass[11pt,a4paper]{article}
\usepackage{dataterra-report}
\usepackage{subcaption}
\usepackage{float}
\usepackage{hyperref}
\usepackage{xcolor}
\usepackage{verbatim} 
\definecolor{airforceblue}{rgb}{0.36, 0.54, 0.66}
\definecolor{ForestGreen}{rgb}{0.13, 0.55, 0.13}
\definecolor{orange}{rgb}{1.0, 0.5, 0.0}
\definecolor{amethyst}{rgb}{0.6, 0.4, 0.8}
\definecolor{awesome}{rgb}{1.0, 0.13, 0.32}
\definecolor{chromeyellow}{rgb}{1.0, 0.65, 0.0}
\hypersetup{colorlinks=true, linkcolor=amethyst, citecolor= ForestGreen, urlcolor=blue } 

% let long property names break after the colon, and give paragraphs a little room
\let\dtprop\prop
\renewcommand{\prop}[1]{\dtprop{\colonbreak#1:\relax}}
\def\colonbreak#1:#2\relax{#1\ifx\relax#2\relax\else:\allowbreak\colonbreak#2\relax\fi}
\setlength{\emergencystretch}{2.5em}

\reportsubtitle{Two observations on the Metadata Guideline and the EarthPortal schema}

\begin{document}
\makereporttitlepage

% =====================================================================
\section{Introduction}

In this short report, we highlight a few observations we made while developing our pipeline on
the state of the art of the semantic artefacts hosted on EarthPortal. Even though we will discuss
that pipeline at the end of October, please \textit{think about our two concerns} related to the Metadata
Guideline:
\begin{enumerate}
  \item \textit{some properties of the guideline (MOD~3) do not have the same name in the metadata schema of
    EarthPortal (MOD~1.4), and five of them are not linked to it at all};
  \item \textit{some fields of the metadata schema are expected by the portal, and most of them by
    O'FAIRe, but they are not in the guideline.}
\end{enumerate}

The guideline is taken in its machine-readable form, the cheatsheet of Report~1 (99 properties).
The metadata schema is the one the portal publishes through its API\footnote{%
\url{https://data.earthportal.eu/submission_metadata} (110 fields) and
\url{https://data.earthportal.eu/ontology_metadata} (20 fields).}. 

% =====================================================================
\section{Development}

\subsection{Observation 1: some properties are stored under another name}

The cheatsheet stores 99 properties: 20 are Mandatory, 58 Recommended and 21 Optional. In
notebook~04\footnote{\url{https://github.com/EarthPortal/SA_Management/tree/main/earthportal_sa_audit}}, 
each of them is linked to a field of the metadata schema with one of three rules
(Table~\ref{tab:rules}).

\begin{center}\small\renewcommand{\arraystretch}{1.2}
\begin{tabular}{@{}l>{\raggedright\arraybackslash}p{10.6cm}r@{}}
\toprule
\textbf{Rule} & \textbf{Meaning} & \textbf{Properties} \\
\midrule
same    & The portal has a field with exactly the same name. & 72 \\
mapping & The portal has a field with another name, and this field lists the property of the
          cheatsheet among its synonyms (\prop{metadataMappings}). & 21 \\
manual  & No rule above works; we wrote the link by hand in \prop{reference/manual\_links.csv}. & 6 \\
\midrule
        & \textbf{Total} & \textbf{99} \\
\bottomrule
\end{tabular}
\captionof{table}{How the 99 properties of the cheatsheet are linked to the EarthPortal schema.}
\label{tab:rules}
\end{center}

Figure~\ref{fig:synonyms} shows how a synonym is declared. The field \prop{omv:description} (Figure~\ref{fig:description})
lists six properties in \prop{metadataMappings}; so, when an ontology file describes itself with
\prop{dcterms:description} or \prop{schema:description}, the portal stores the value in
\prop{omv:description}. For Description, the cheatsheet already uses \prop{omv:description}, so
it is a \emph{same} case; but Keywords works exactly in this way: the cheatsheet says
\prop{dcat:keyword}, and the portal field \prop{omv:keywords} lists \prop{dcat:keyword} among its
synonyms (Figure~\ref{fig:keywords}), so it is a \emph{mapping} case.


\begin{figure}[H]
  \centering
  % First Subfigure: Description
  \begin{subfigure}[b]{0.45\textwidth}
    \centering
    \includegraphics[width=\textwidth]{omvDescription_s.png}
    \caption{\texttt{omv:description} (same case)}
    \label{fig:description}
  \end{subfigure}
  \hfill % Adds horizontal stretching space between plots
  % Second Subfigure: Keywords
  \begin{subfigure}[b]{0.50\textwidth}
    \centering
    \includegraphics[width=\textwidth]{omvKeywords.png}
    \caption{\texttt{omv:keywords} (mapping case)}
    \label{fig:keywords}
  \end{subfigure}
  \caption{Equivalence between the properties of the cheatsheet and the fields of the EarthPortal schema. The list
  \prop{metadataMappings} gives the properties the portal accepts as synonyms.} 
  \label{fig:synonyms}
\end{figure}



Figure~\ref{fig:26} lists the 26 properties of the cheatsheet that the portal stores under
another name. Twenty-one of them are \emph{mapping} cases: the portal already knows the name used
by the cheatsheet, so nothing has to change for them. The other five are \emph{manual} cases (the
sixth manual case, Categories, has the same name as the portal field; it is manual only because
\prop{omv:hasDomain} exists on both the ontology record and the submission).

\begin{figure}[H]
  \centering
  \includegraphics[width=\linewidth]{Cheatsheet_vs_earthportal_metadata_schema.png}
  \caption{Output of notebook~04: the 26 properties of the cheatsheet that the portal stores under
  another name, with the rule that links them. For more details on the manual cases, click 
  \href{https://github.com/EarthPortal/SA_Management/blob/main/earthportal_sa_audit/reference/manual_links.csv}{here}.}
  \label{fig:26}
\end{figure}

Table~\ref{tab:five} gives these five properties. For each of them, the portal field has no
synonym that matches the cheatsheet. So, if an author follows the guideline and writes, for
example, \prop{mod:semanticArtefactRelation} in the ontology file, we expect the portal not to
recognise it. To understand where each name comes
from, we also looked at the current version of MOD (3.2.5)\footnote{File \prop{mod.ttl} of
\url{https://github.com/FAIR-IMPACT/MOD}, branch \prop{main}, read on 14 October 2026.}.
% MOD 3.0: https://github.com/FAIR-IMPACT/MOD/blob/main/mod.ttl

\begin{center}\footnotesize\renewcommand{\arraystretch}{1.3}\let\dtprop\texttt
\begin{tabular}{@{}l>{\raggedright\arraybackslash}p{2.0cm}>{\raggedright\arraybackslash}p{3.75cm}>{\raggedright\arraybackslash}p{3.55cm}>{\raggedright\arraybackslash}p{4.45cm}@{}}
\toprule
\textbf{Id} & \textbf{Entry} & \textbf{Cheatsheet} & \textbf{EarthPortal} & \textbf{What we found} \\
\midrule
1.1  & Acronym & \prop{mod:acronym} & \prop{omv:acronym} (ontology record) &
  MOD 3.2 declares \prop{mod:acronym}; the portal field has no synonym. \\
1.9  & Version IRI & \prop{omv:versionIRI} & \prop{owl:versionIRI} &
  MOD 3.2 uses \prop{owl:versionIRI}; we did not find \prop{omv:versionIRI} in MOD. Here, the
  guideline may be the one to correct. \\
10.1 & This ontology is a view of & \prop{dct:isPartOf} & \prop{bpm:viewOf} (ontology record) &
  The portal has no field for \prop{dct:isPartOf}; it has the inverse,
  \prop{dct:hasPart} (``Has part (has views)''). \\
10.5 & Generally related to & \prop{mod:semanticArtefact\allowbreak Relation} & \prop{door:ontologyRelatedTo} &
  MOD 3.2 declares both properties equivalent, but the portal field does not list the MOD name. \\
10.10 & Disparate modelling with & \prop{mod:hasDisparateModelling} & \prop{door:hasDisparateModelling} &
  The portal lists \prop{mod:hasDisparateModellingWith}, while \prop{mod.ttl} declares
  \prop{mod:hasDisparateModelling}: one spelling has to be aligned. \\
\bottomrule
\end{tabular}
\captionof{table}{The five properties of the cheatsheet that are not linked to the EarthPortal
schema, neither by name nor by a synonym.}
\label{tab:five}
\end{center}

\begin{keypoint}
Before the publication of this guideline, it could be better for the development team to update
the code and declare new mappings for these five properties, or to agree with us on the name the
guideline should use. The 21 mapped properties need nothing.
\end{keypoint}

% ---------------------------------------------------------------------
\subsection{Observation 2: some fields expected by the portal are not in the guideline}

We then looked at the other direction: the fields of the metadata schema that no property of the
guideline points to. Of the 130 fields (110 on the submission, 20 on the ontology record), 31 are
not used by the guideline:
\begin{itemize}
  \item 21 are managed by the portal itself: they have no label, their category is ``no'', and they
    are not extracted from the ontology file (for example \prop{submissionId},
    \prop{submissionStatus}, \prop{uploadFilePath}, \prop{metrics}, \prop{notes}, \prop{reviews},
    \prop{views}). They do not belong in a guideline for authors;
  \item 10 have a label in the submission form and are extracted from the ontology file. They are
    the candidates of Table~\ref{tab:candidates}.
\end{itemize}

The last two columns of Table~\ref{tab:candidates} say whether O'FAIRe reads the field (with the
question and its points), and in how many of the 68 artefacts (views and metadata-only records
excluded) the latest submission has a value today.

\begin{center}\small\renewcommand{\arraystretch}{1.25}
\begin{tabular}{@{}l>{\raggedright\arraybackslash}p{4.0cm}>{\raggedright\arraybackslash}p{1.7cm}>{\raggedright\arraybackslash}p{3.2cm}r@{}}
\toprule
\textbf{Portal field} & \textbf{Label in the portal} & \textbf{Category} & \textbf{O'FAIRe question (points)} & \textbf{Filled} \\
\midrule
\prop{omv:version}            & Version information               & general & R1.2Q4 (4)          & 56 / 68 \\
\prop{owl:deprecated}         & Deprecated                        & general & A2Q4 (4)            & 55 / 68 \\
\prop{bpm:prefLabelProperty}  & Object preferred label property   & object  & I3Q3 (7), R1Q1 (5)  & 55 / 68 \\
\prop{bpm:synonymProperty}    & Object synonym property           & object  & I3Q3 (7), R1Q1 (5)  & 55 / 68 \\
\prop{bpm:definitionProperty} & Object definition property        & object  & I3Q3 (7), R1Q1 (5)  & 55 / 68 \\
\prop{bpm:authorProperty}     & Object author property            & object  & I3Q3 (7), R1Q1 (5)  & 55 / 68 \\
\prop{bpm:obsoleteProperty}   & Object obsolete property          & object  & I3Q3 (7), R1Q1 (5)  & 55 / 68 \\
\prop{bpm:hierarchyProperty}  & Hierarchy property                & object  & I3Q3 (7), R1Q2 (3)  & 57 / 68 \\
\prop{bpm:createdProperty}    & Object creation date property     & object  & no                       & 55 / 68 \\
\prop{bpm:modifiedProperty}   & Object modification date property & object  & no                       & 55 / 68 \\
\bottomrule
\end{tabular}
\captionof{table}{Fields of the EarthPortal schema that could be added to the guideline.
``object'' stands for the category ``object description properties''. In O'FAIRe, MOD~1.4 calls
the first field \prop{owl:versionInfo} and the object description fields \prop{mod:...Property}.}
\label{tab:candidates}
\end{center}

Three remarks on this table:
\begin{enumerate}
  \item Eight of the ten fields bring points in O'FAIRe, yet an author who follows the guideline
    has no way to know that they exist. The guideline already has Version IRI (1.9) and Prior
    version (10.2), but not the version number itself; and it already has Root of obsolete branch
    (11.1, \prop{bpm:obsoleteParent}), which describes the content in the same way as the object
    description fields.
  \item The values look like default values rather than choices of the authors. In the 55 filled
    submissions, each object description field has exactly one value: for instance,
    \prop{prefLabelProperty} is \prop{skos:prefLabel} everywhere, including in 23 OWL artefacts,
    and \prop{deprecated} is \prop{false} everywhere. We have not checked where these defaults come
    from.
  \item Some of these values are not in the list of values the portal itself proposes
    (\prop{enforcedValues}). For 22 OWL artefacts, \prop{hierarchyProperty} is
    \prop{owl:subClassOf}, while the list proposes \prop{rdfs:subClassOf}; and
    \prop{createdProperty} and \prop{modifiedProperty} are \prop{dc:created} and
    \prop{dc:modified} (Dublin Core Elements 1.1), while the list proposes \prop{dcterms:created}
    and \prop{dcterms:modified}. 
\end{enumerate}

\begin{keypoint}
The ten fields of Table~\ref{tab:candidates} could be added to the guideline, at least the eight
that O'FAIRe reads. If they are added, the guideline could also give the allowed values, so that
authors replace the defaults with the properties their artefact really uses.
\end{keypoint}

% =====================================================================
\section{Conclusion}

While I am still completing and building some intermediary steps within my pipeline (five Jupyter
notebooks), I would appreciate it if you could check my two statements and provide some feedback:
\begin{enumerate}
  \item five properties of the guideline are not linked to the EarthPortal schema
    (Table~\ref{tab:five}), and the development team could declare new mappings for them before
    the guideline is published;
  \item ten fields of the schema, eight of them read by O'FAIRe, could be added to the guideline
    (Table~\ref{tab:candidates}).
\end{enumerate}
Everything can be reproduced with notebook~04 of our repository\footnote{%
\url{https://github.com/EarthPortal/SA_Management}, folder \prop{earthportal\_sa\_audit}.}.

\end{document}
